\section*{CHAPTER 1. INTRODUCTION AND LITERATURE REVIEW}
\addcontentsline{toc}{section}{\numberline{}CHAPTER 1. INTRODUCTION AND LITERATURE REVIEW}

\setcounter{section}{1}
\setcounter{subsection}{0}
\setcounter{figure}{0}
\setcounter{table}{0}

\subsection{Macro View: The Evolution of Smart Grid and Multi-Microgrids}
The global energy landscape is undergoing a profound paradigm shift, DERs \cite{Muhtadi2021, Dagar2021}. The massive penetration of renewables has catalyzed the transformation of passive end-users into active prosumers. Consequently, this evolution introduces unprecedented challenges in two-way energy coordination and grid management \cite{Ahmad2023}.

To manage the escalating complexity of this decentralized environment, MG has emerged as the fundamental cellular building block of the modern Smart Grid \cite{She2022}. Functioning as a localized energy system, a microgrid possesses the critical capability to operate seamlessly in both grid-connected and isolated (islanded) modes \cite{Trivedi2022}. Furthermore, BESS have become an indispensable component within these microgrids, serving as a vital buffer to absorb the inherent stochasticity of renewable generation \cite{Choudhury2022}.

In parallel with physical infrastructure advancements, the economic mechanisms governing energy exchange are fundamentally evolving. Rigid incentive structures, FIT, P2P energy trading frameworks \cite{Zhang2018_P2P, Karami2021}. P2P trading facilitates direct active power transactions among localized prosumers \cite{Ali2022}, thereby maximizing community social welfare and significantly shortening investment payback periods \cite{Park2024, Tushar2020_P2P}. Ultimately, this cultivates a high degree of regional energy autonomy \cite{Zia2020}. It is crucial to emphasize that within this framework, P2P markets are strictly restricted to the trading of active power ($P$). Reactive power ($Q$), conversely, is not subjected to financial transactions but is entirely generated and managed internally within each microgrid to support local voltage stability.

More recently, the strategic focus of microgrid and P2P implementation has necessitated a critical pivot from pure economic efficiency toward grid resilience. Driven by the escalating frequency of extreme weather events, networked microgrids are increasingly conceptualized as critical safety nets \cite{Li2017Networked, Hu2026Region}. During severe grid disturbances, these systems must provide self-healing capabilities and facilitate emergency energy routing to safeguard critical loads \cite{Chen2020}. Consequently, the localized P2P trading platform transcends its conventional financial role, functioning fundamentally as the underlying physical infrastructure that enables such emergency energy routing during autonomous islanding operations. This ensures that operational continuity and power system resilience are robustly maintained even under extreme adverse conditions \cite{Bhusal2020}.

\subsection{The Physical Layer: Power Flow Models and Network Constraints}
Early literature frequently models distribution grids and peer-to-peer markets using a lumped-node (or copper-plate) abstraction. This oversimplified approach aggregates all local generations and loads into a single node, operating under the assumption of infinite transmission capacity and ignoring the physical network topology \cite{Guerrero2021_Network}. By disregarding line impedances and thermal limits, the copper-plate model assumes that active power can be freely routed between any two prosumers without encountering physical constraints. Consequently, such models obscure underlying voltage violations and often trigger severe localized congestion when deployed in real-world scenarios \cite{Guerrero2018_Decentralized}. As demonstrated by recent studies, relying on lumped-node aggregations creates a misleading perception of grid flexibility, which inevitably fails during high-stress operating conditions \cite{lannoye2012evaluation, javed2024analyzing}.

To incorporate network constraints, subsequent research adopted linearized models like DC Optimal Power Flow (DC-OPF) or LinDistFlow. These models operate by neglecting branch losses and transverse voltage drops, essentially dropping the non-linear terms from the branch flow equations \cite{paudel2020peer}. While computationally efficient for bulk transmission systems, this mathematical simplification introduces significant errors in distribution networks characterized by high R/X ratios \cite{Patari2021}. The inherent physical characteristics of radial distribution grids dictate that active and reactive power (P-Q-V) are strongly coupled \cite{Hashmi2020}. Any rapid fluctuation in active power ($P$)—such as those induced by aggressive decentralized P2P trading—immediately induces substantial variations in nodal voltage magnitudes. Therefore, decoupling the economic dispatch of active power from the physical management of reactive power leads to severely underestimated voltage drops \cite{Pareek2020}. Under extreme weather events or emergency islanded modes, relying on such linear approximations frequently leads to voltage collapse \cite{Polat2025, eajal2018existence}.

AC-OPF problem, SOCP relaxation has emerged as a robust framework. The core conflict in AC-OPF arises from the non-linear coupling of voltage magnitudes and phase angles, which inherently creates a non-convex feasible region. SOCP deconstructs this conflict by defining auxiliary variables to represent voltage products---such as mapping $W_{ij} = V_{i} V_{j} \cos(\theta_{ij})$---and applying exact convexification. Specifically, it relaxes the rigid non-linear equality constraints into a convex inequality cone, yielding relationships such as $W_{ii} W_{jj} \geq |W_{ij}|^2$ \cite{farivar2013branch, Patari2021}. By transforming the non-convex domain into a convex space, this mathematical mechanism guarantees exactness and global optimality specifically under radial network topologies, while accelerating computational speed and inherently preserving the strict AC physics of the grid \cite{Low2014_SOCP}\footnote{It is critical to acknowledge that the strict exactness of the SOCP formulation is theoretically maintained because the overarching objective function incorporates a Value of Lost Load (VoLL) penalty component. This mechanism, combined with the radial topology of the microgrid, mathematically penalizes load shedding, thereby implicitly forcing the system to tighten the SOC convergence boundary (SOC relaxation tightness) and ensuring physically realizable voltage bounds.}.

More importantly, unlike linear approximations that ignore losses, the SOCP formulation accurately captures the exact quadratic power flow equations, longitudinal branch losses, and strict transverse voltage limits essential for resilient operation \cite{Low2014_SOCP, Srivastava2023}. This precise physical representation becomes critical during emergency islanding conditions, where the microgrid operates independently and is susceptible to voltage collapse. The SOCP framework effectively encapsulates the unique islanding physics of microgrids by tightly coupling the active power routing with the underlying voltage profile. By respecting the true non-linear bounds of the network, the optimization ensures that the derived market clearing solutions remain strictly physically feasible under extreme conditions, avoiding the misleading perception of grid flexibility.

Within this exact AC framework, the critical role of reactive power ($Q$) is explicitly managed. Rather than commoditizing reactive power into the financial Peer-to-Peer market, the SOCP model mathematically ensures that local distributed energy resources provide autonomous $Q$ generation. Specifically, DERs autonomously generate $Q$ as an ancillary service utilizing their remaining inverter capacity ($S_{max}^2 \ge P^2 + Q^2$) \cite{Montoya2021Mixed}, thus preserving the active power ($P$) trading limits \cite{Angle2020SOCP}. This localized reactive support is dynamically dispatched to neutralize the voltage fluctuations induced by large-scale active power ($P$) trades. As demonstrated by Garrido-Arevalo et al., integrating this mechanism guarantees that local grid security is maintained, proving effective in supporting voltage stability without requiring external financial transactions \cite{Garrido_Arevalo_2024}. Furthermore, Men et al. explicitly established that deploying such SOCP-based autonomous reactive compensation enhances overall power system resilience, ensuring uninterrupted operation during severe disturbances and autonomous islanded transitions \cite{men2023secondary}.

\subsection{The Spatial Layer: Decentralization and Privacy}
Traditional energy management systems depend on a centralized architecture to coordinate dispatch and execute global optimization. However, EMS to aggregate large volumes of high-resolution data, including local capacity bounds and operational states. This extensive data requirement creates a severe computational and communication bottleneck that degrades real-time dispatch efficiency \cite{Chen2021}. Furthermore, P2P market, prosumers' underlying load profiles and financial bids are strictly confidential. A centralized framework inherently compromises this privacy by mandating full parameter exposure, thereby undermining market trust and disincentivizing participation \cite{Liu2025}.
Beyond computational and economic limitations, centralized coordination introduces a structural vulnerability that compromises grid resilience during extreme weather events. Operating a multi-microgrid network through a single computational hub establishes a critical single point of failure. If the central communication link is severed during a physical disaster, the entire multi-microgrid system loses coordination. This loss of centralized control triggers cascading failures across the network, disrupting the autonomous islanded resilience required to sustain critical loads during emergencies \cite{Panteli2016_Resilience, Zuo2022}.

To mitigate these vulnerabilities, distributed optimization frameworks, ADMM, have been widely adopted. ADMM establishes a flat decentralized structure that mathematically decomposes the global optimization problem \cite{Zhao2022}. This mechanism protects data privacy, as participants only exchange boundary coupling variables rather than internal parameters \cite{erseghe2014distributed}. Moreover, this flat consensus-based approach ensures resilience against centralized failures \cite{Zhang2019_ADMM}. Concurrently, DLMP \cite{Ullah2022, Nasiri2023}. These pricing signals accurately reflect local marginal costs and network congestion states \cite{Xia2024}. Consequently, ADMM has proven effective for linearized, flat market clearings \cite{Mishra2024}.

However, modern power systems are inherently hierarchical rather than flat. A distribution network operator typically coordinates multiple autonomous microgrids, forming a strict top-down operational topology. When ADMM is applied to solve such multi-level optimizations, particularly those involving exact SOCP formulations, its mathematical foundation exposes critical limitations. Fundamentally, ADMM enforces ``flat horizontal consensus'', which causes severe oscillation and algorithmic divergence in non-linear SOCP multi-level grids \cite{Rajaei2021}. This inefficiency creates substantial communication overhead, rendering the method inadequate for dynamic, multi-level grid coordination \cite{Sun_2022}.

To resolve this structural mismatch, ATC provides a mathematically rigorous top-down coordination framework. Mathematically, ATC achieves this decoupling by formulating an Augmented Lagrangian penalty function to coordinate boundary coupling variables between hierarchical levels. Crucially, convexifying the physical grid via SOCP is a strict prerequisite condition for this process; without a globally convex feasible region at the physical layer, the ATC algorithm cannot theoretically guarantee convergence. Consequently, by enforcing a ``vertical master-slave consensus'' over these convexified domains, ATC securely decouples the hierarchical regions and ensures fast convergence \cite{mohammadi2020accelerated}. In this paradigm, the upper level dictates optimal ``Target'' variables, while the lower levels compute and return corresponding ``Response'' variables \cite{Zhao2022_ATC, Kong2020_ATC}. This strict target-response mechanism decouples the hierarchical optimization loops, allowing each layer to solve its mathematical problem independently without violating data privacy. Importantly, DLMP, preserving economic pricing signals within a top-down structure. By eliminating the necessity for global flat consensus, ATC significantly reduces communication overhead and guarantees fast, stable convergence even when managing complex, non-linear SOCP models.

\subsection{The Temporal Layer: Temporal Dynamics and Grid Resilience}
Traditional operational strategies predominantly rely on static optimization or pure day-ahead dispatch mechanisms governed by a deterministic mathematical formulation. This approach calculates an operational trajectory predicated entirely on a single day-ahead snapshot, denoted as $\hat{x}$, rendering the control system fundamentally rigid \cite{Liu2020Robust, Huynh2024Water}. Because it lacks a continuous update loop, the mathematical model cannot dynamically adapt when real-time grid conditions inevitably deviate from the initial snapshot $\hat{x}$ \cite{zhang2021multi, wu2021model}. Consequently, the pre-calculated dispatch bounds are rapidly breached when the power system encounters sudden physical disturbances, such as emergency islanding transitions. Enforcing this rigid schedule under such conditions causes immediate voltage deviations and capacity constraint violations, which can trigger severe cascading failures across the network.

To address the mathematical limitations of deterministic static models, SO has been widely explored. SO formulates the dispatch problem by optimizing the expected value of an objective function, typically expressed as $E[f(x, \xi)]$, PDFs. While this approach mitigates minor forecasting errors under normal operations, it exhibits critical theoretical vulnerabilities when applied to grid resilience. Specifically, relying on established PDFs is often insufficient to capture High-Impact, HILP events or fat-tail risks. Because these extreme physical disruptions lack sufficient historical data, synthesizing accurate statistical distributions is infeasible, rendering the SO approach inadequate for safeguarding critical infrastructure during severe blackouts.

Alternatively, RO circumvents the reliance on precise statistical distributions by framing the control problem within a strict min-max formulation, represented as $\min \max f(x, u)$. This framework evaluates the worst-case scenario over a predefined uncertainty set $U$, guaranteeing operational feasibility under extreme conditions. However, mathematically preparing for the absolute worst-case creates extreme over-conservatism, forcing the microgrid to maintain substantial idle energy reserves and substantially degrading economic efficiency. Furthermore, constructing a rigorous two-stage RO model for complex microgrid networks inevitably leads to highly non-linear min-max-min mathematical structures. These structures are computationally intractable and cannot be solved efficiently enough for real-time applications. Crucially, despite their probabilistic or worst-case bounds, both SO and RO remain open-loop strategies that fundamentally lack a real-time temporal feedback loop to continuously correct dispatch decisions as extreme events unfold.

Overcoming the structural vulnerabilities of both static and conservative paradigms, MPC introduces a dynamic, multi-temporal optimization framework centered on a rolling horizon mechanism. Mathematically, the MPC control loop solves an open-loop optimal dispatch problem over a predefined prediction horizon $H$. However, the core mechanism of this framework lies in its state-space feedback loop: the controller executes only the optimal control action for the immediate first time step, denoted as $u_k^*$ \cite{Elkazaz2019Energy}. Following this execution, the algorithm continuously updates the actual system state $x_{k+1}$ based on real-time sensor measurements and shifts the prediction window forward. This recursive mathematical adaptability effectively neutralizes both forecasting errors and sudden physical disturbances by continuously recalculating the optimal trajectory in real-time \cite{Parisio2014_MPC}. Furthermore, integrating real-time MPC with iterative ATC is computationally feasible because it leverages a ``warm-start'' mechanism---reusing the previous time-step's optimum---which drastically reduces the required ATC iterations \cite{zheng2023distributed}.

Within this dynamic rolling framework, the MPC objective function is explicitly formulated to enhance grid survivability during extreme events by incorporating a Value of Lost Load (VoLL) penalty, represented as $C_{VoLL}$. By mathematically weighting this penalty according to load criticality, the MPC implements a physical execution strategy known as graceful degradation. As the microgrid's battery State of Charge (SoC) inevitably depletes during extended blackouts, the algorithm utilizes the predictive foresight of the sliding window to evaluate future capacity limits. It then proactively sheds low-priority demands long before the battery is fully exhausted \cite{zhou2025networked}. This proactive energy routing mathematically guarantees the survival of critical infrastructure throughout prolonged disruptions, thereby minimizing the total VoLL penalty and maximizing physical grid resilience. It must be explicitly stated, however, that this MPC model manages operations strictly at the energy scheduling level (Tertiary and Secondary Control). The instantaneous inter-step dynamic stability, occurring within the millisecond timeframe, is assumed to be independently governed by the underlying Primary Control layer.

\subsection{Research Gap, Motivation, and Major Contributions}

\subsubsection{The Master Gap and Motivation: Overcoming Fragmented Paradigms}
Despite significant advancements in microgrid operation, a critical review of the existing literature reveals that current methodologies overwhelmingly address physical constraints, spatial coordination, and temporal dynamics in isolation. Specifically, three major limitations persist:

\begin{itemize}
    \item \textbf{Neglect of Non-linear Physics:} Distributed optimization techniques, notably ADMM, are frequently applied alongside linearized power flow models. While computationally efficient, this combination inherently ignores the critical coupling of active and reactive power, completely neglecting reactive power ($Q$) and strict voltage limits, which leads to voltage collapse during emergency operations.
    \item \textbf{Privacy and Structural Vulnerabilities:} SOCP EMS. P2P trading.
    \item \textbf{Static Temporal Rigidity:} ATC are often deployed within static, day-ahead environments. Because these static models lack a real-time temporal feedback loop, they are mathematically rigid and are highly vulnerable to the unpredictable nature of extreme High-Impact, HILP events.
\end{itemize}

The master gap in contemporary research, therefore, SOCP, ATC, MPC. The primary motivation of this thesis is to bridge this gap by transitioning P2P energy coordination from pure financial market trading toward "Cooperative Resilience", wherein localized energy exchange serves as the physical backbone for surviving severe grid disturbances.

\subsubsection{Major Contributions}
To address these critical limitations, this research proposes a comprehensive framework that unifies the physical, spatial, and temporal dimensions of multi-microgrid operations. The major contributions of this thesis are defined by four foundational pillars:

\begin{enumerate}
    \item \textbf{SOCP:} AC-OPF model where P2P trading is strictly limited to the exchange of active power ($P$). Concurrently, reactive power ($Q$) DERs, ensuring voltage stability without the need for financial commoditization.
    \item \textbf{ATC:} ATC framework to securely coordinate interconnected multi-microgrid networks. This hierarchical mechanism ensures strict data privacy and structural resilience, effectively resolving the slow convergence and oscillation issues inherent to ADMM when applied to complex, non-linear physical systems.
    \item \textbf{MPC:} MPC strategy, embedded with a Value of Lost Load (VoLL) penalty. This integration transforms the traditional static dispatch model into a dynamic control framework, mathematically capable of executing "Graceful Degradation" to proactively safeguard critical loads and sustain critical operations during prolonged outages.
    \item \textbf{Integration and Validation:} Rigorously simulating the unified SOCP-ATC-MPC architecture on a complex, hierarchical multi-microgrid testbed under extreme physical disturbances. The comprehensive validation demonstrates the effectiveness of the proposed tri-layer framework over conventional, isolated methodologies in maintaining operational continuity.
\end{enumerate}
